% !TEX program = pdflatex
\documentclass[11pt,a4paper]{article}

\usepackage[margin=1in]{geometry}
\usepackage{parskip}

\usepackage{amsmath,amssymb,amsthm}
\usepackage{xcolor}

\theoremstyle{definition}
\newtheorem{definition}{Definition}
\newtheorem{example}{Example}

\theoremstyle{plain}
\newtheorem{theorem}{Theorem}
\newtheorem{proposition}{Proposition}
\newtheorem{lemma}{Lemma}

\usepackage{mathtools,longtable,booktabs,array}
\setlength{\emergencystretch}{2em}
\allowdisplaybreaks[1]
\raggedbottom
\begin{document}


\begin{center}
    {\LARGE \textbf{Mathematical Methods for Physics--Lecture 14}} \\[0.6em]
    {\large The Delta Function and Fourier Transform Properties}
\end{center}

\vspace{-0.5em}
\noindent\rule{\textwidth}{0.5pt}
\vspace{0.3em}

\begin{center}
    Instructor: Zhaolong Han, SOCE, Shanghai Jiao Tong University;
\end{center}

\vspace{-0.6em}
\noindent Email: han.arkey@sjtu.edu.cn

\vspace{0.4em}
\noindent\rule{\textwidth}{0.5pt}

\vspace{0.8em}

\section{The Delta Function and Generalized Fourier Transforms}

\subsection{The Delta Function and Its Rectangular-Pulse Approximation}

The \(\delta\) function is also called the unit impulse. It is a distribution,
formally represented by
\[
\delta(x)=
\begin{cases}
0,&x\ne0,\\
\infty,&x=0,
\end{cases}
\qquad
\int_{-\infty}^{+\infty}\delta(x)\,\mathrm dx=1.
\]

\begin{definition}[Rectangular-pulse representation]
\[
\delta(x)=\lim_{\tau\to0}\delta_\tau(x),
\]

where \(\tau>0\) and
\[
\delta_\tau(x)=
\begin{cases}
\dfrac1{\tau},&|x|\le\dfrac{\tau}{2},\\[4pt]
0,&|x|>\dfrac{\tau}{2}.
\end{cases}
\]

\end{definition}

The limit is understood in the distributional sense. The rectangular pulse is
centered at the origin, has width \(\tau\) and height \(1/\tau\), and
has unit area. As \(\tau\to0\), it becomes narrower and taller. A vertical arrow at the origin labeled ``1'' represents the unit impulse.

\subsection{The Delta Function as a Functional}

\begin{definition}[The Dirac delta distribution]

The delta function \(\delta(t)\) is a functional satisfying
\[
\int_{-\infty}^{+\infty}
\delta(t-t_0)\varphi(t)\,\mathrm dt
=\varphi(t_0),
\qquad \forall\varphi\in C(\mathbb R).
\]

\end{definition}

Its derivative satisfies
\[
\int_{-\infty}^{+\infty}
\delta'(t-t_0)\varphi(t)\,\mathrm dt
=-\varphi'(t_0),
\qquad \forall\varphi\in C^1(\mathbb R).
\]

Its \(n\)th derivative satisfies
\[
\int_{-\infty}^{+\infty}
\delta^{(n)}(t-t_0)\varphi(t)\,\mathrm dt
=(-1)^n\varphi^{(n)}(t_0),
\qquad \forall\varphi\in C^n(\mathbb R).
\]

\subsection{Fourier Transform and Integral Representation of the Delta Function}

\textbf{Example 1.} The Fourier transform of \(\delta(t)\) is understood
as a generalized Fourier transform.
\[
\mathcal F[\delta(t)]
=\int_{-\infty}^{+\infty}\delta(t)e^{-i\omega t}\,\mathrm dt
=\left.e^{-i\omega t}\right|_{t=0}
=1.
\]

Thus \(\{\delta(t),1\}\) is a Fourier transform pair. By the inverse
transform,
\[
\delta(t)=\frac1{2\pi}
\int_{-\infty}^{+\infty}1\cdot e^{i\omega t}\,\mathrm d\omega.
\]

Consequently,
\[
\int_{-\infty}^{+\infty}e^{\pm i\omega t}\,\mathrm dt
=2\pi\delta(\omega).
\tag{*}
\]

\subsection{Transforms of Complex Exponentials, Cosines, and Sines}

\textbf{Example.} Prove that
\[
e^{i\omega_0t}
\ \xleftrightarrow{\mathcal F}\
2\pi\delta(\omega-\omega_0).
\]

Using identity \((*)\),
\[
\begin{aligned}
\mathcal F[e^{i\omega_0t}]
&=\int_{-\infty}^{+\infty}
e^{i\omega_0t}e^{-i\omega t}\,\mathrm dt\\
&=\int_{-\infty}^{+\infty}
e^{-i(\omega-\omega_0)t}\,\mathrm dt\\
&=2\pi\delta(\omega-\omega_0).
\end{aligned}
\]

Similarly,
\[
\begin{aligned}
\mathcal F[\cos\omega_0t]
&=\mathcal F\!\left[
\frac{e^{i\omega_0t}+e^{-i\omega_0t}}2
\right]\\
&=\frac12\mathcal F[e^{i\omega_0t}]
+\frac12\mathcal F[e^{-i\omega_0t}]\\
&=\frac12\cdot2\pi\delta(\omega-\omega_0)
+\frac12\cdot2\pi\delta(\omega+\omega_0)\\
&=\pi\bigl[\delta(\omega+\omega_0)+\delta(\omega-\omega_0)\bigr].
\end{aligned}
\]
\[
\mathcal F[\sin\omega_0t]
=i\pi\bigl[\delta(\omega+\omega_0)-\delta(\omega-\omega_0)\bigr].
\]

\subsection{Transform of the Unit Step Function}
\[
u(t)=
\begin{cases}
1,&t>0,\\
0,&t<0.
\end{cases}
\]

Its Fourier transform is
\[
u(t)
\ \xleftrightarrow{\mathcal F}\
\frac1{i\omega}+\pi\delta(\omega).
\]

Here \(1/(i\omega)\) is interpreted as a Cauchy principal-value distribution.
The forward transform can formally be written as
\[
\mathcal F[u(t)]
=\int_{-\infty}^{+\infty}u(t)e^{-i\omega t}\,\mathrm dt
=\int_0^{+\infty}e^{-i\omega t}\,\mathrm dt.
\]

Verify the result using the inverse transform.
\[
\begin{aligned}
\mathcal F^{-1}\!\left[
\frac1{i\omega}+\pi\delta(\omega)
\right]
&=\frac1{2\pi}\int_{-\infty}^{+\infty}
\left[\frac1{i\omega}+\pi\delta(\omega)\right]
e^{i\omega t}\,\mathrm d\omega\\
&=\frac1{2\pi}\left[
\int_{-\infty}^{+\infty}
\frac{e^{i\omega t}}{i\omega}\,\mathrm d\omega
+\int_{-\infty}^{+\infty}
\pi\delta(\omega)e^{i\omega t}\,\mathrm d\omega
\right]\\
&=\frac1{2\pi}\left[
\int_{-\infty}^{+\infty}
\frac{e^{i\omega t}}{i\omega}\,\mathrm d\omega+\pi
\right].
\end{aligned}
\]

where
\[
\begin{aligned}
\int_{-\infty}^{+\infty}
\frac{e^{i\omega t}}{i\omega}\,\mathrm d\omega
&=\int_{-\infty}^{+\infty}
\frac{\cos\omega t+i\sin\omega t}{i\omega}\,\mathrm d\omega\\
&=0+\int_{-\infty}^{+\infty}
\frac{\sin\omega t}{\omega}\,\mathrm d\omega\\
&=
\begin{cases}
\pi,&t>0,\\
-\pi,&t<0.
\end{cases}
\end{aligned}
\]

From
\[
\int_{-\infty}^{+\infty}\frac{\sin x}{x}\,\mathrm dx=\pi
\]

we obtain
\[
\int_{-\infty}^{+\infty}\frac{\sin kx}{x}\,\mathrm dx
=
\begin{cases}
\pi,&k>0,\\
-\pi,&k<0.
\end{cases}
\]

Substituting into the preceding expression gives
\[
\mathcal F^{-1}\!\left[
\frac1{i\omega}+\pi\delta(\omega)
\right]
=
\begin{cases}
\dfrac1{2\pi}(\pi+\pi)=1,&t>0,\\[4pt]
\dfrac1{2\pi}(-\pi+\pi)=0,&t<0,
\end{cases}
=u(t).
\]

\subsection{A Constant (DC) Signal}

The transform pair for a constant (DC) signal \(E\) is
\[
E\ \xleftrightarrow{\mathcal F}\ 2\pi E\delta(\omega).
\]
\[
\mathcal F[E]
=\int_{-\infty}^{+\infty}Ee^{-i\omega t}\,\mathrm dt
=2\pi E\delta(\omega).
\]

\section{Properties of Fourier Transform Pairs}

We use the notation
\[
f(t)\ \xleftrightarrow{\mathcal F}\ F(\omega),
\qquad
g(t)\ \xleftrightarrow{\mathcal F}\ G(\omega).
\]

The transform convention is
\[
F(\omega)=\int_{-\infty}^{+\infty}f(t)e^{-i\omega t}\,\mathrm dt,
\qquad
f(t)=\frac1{2\pi}\int_{-\infty}^{+\infty}F(\omega)e^{i\omega t}\,\mathrm d\omega.
\]

\subsection{Linearity}
\[
\begin{aligned}
\mathcal F[af(t)+bg(t)]
&=a\mathcal F[f]+b\mathcal F[g]\\
&=aF(\omega)+bG(\omega).
\end{aligned}
\]

Alternatively,
\[
\mathcal F^{-1}[aF(\omega)+bG(\omega)]
=a\mathcal F^{-1}[F]+b\mathcal F^{-1}[G].
\]

\subsubsection*{Example 7.10}
\[
\begin{aligned}
\mathcal F[\sin^2t]
&=\mathcal F\!\left[\frac{1-\cos2t}{2}\right]\\
&=\frac12\left\{\mathcal F[1]-\mathcal F[\cos2t]\right\}\\
&=\frac12\left\{
2\pi\delta(\omega)
-\pi[\delta(\omega+2)+\delta(\omega-2)]
\right\}.
\end{aligned}
\]

\subsubsection*{Example 7.11}

Let
\[
F(\omega)=\pi\delta(\omega+1)-\frac{i}{\omega+1},
\]

Find \(\mathcal F^{-1}[F(\omega)]\).

\textbf{Solution}\par\nopagebreak[4]
\[
\begin{aligned}
\mathcal F^{-1}[F(\omega)]
&=\frac1{2\pi}\int_{-\infty}^{+\infty}
\left[\pi\delta(\omega+1)-\frac{i}{\omega+1}\right]
e^{i\omega t}\,\mathrm d\omega\\
&=\frac1{2\pi}\left[
\int_{-\infty}^{+\infty}
\pi\delta(\omega+1)e^{i\omega t}\,\mathrm d\omega
+\int_{-\infty}^{+\infty}
\frac{-ie^{i\omega t}}{\omega+1}\,\mathrm d\omega
\right]\\
&=\frac1{2\pi}\left[
\pi e^{-it}
+\int_{-\infty}^{+\infty}
\frac{-ie^{i\omega t}}{\omega+1}\,\mathrm d\omega
\right].
\end{aligned}
\]

Let \(\omega+1=\xi\). Then
\[
\begin{aligned}
\int_{-\infty}^{+\infty}
\frac{-ie^{i\omega t}}{\omega+1}\,\mathrm d\omega
&=\int_{-\infty}^{+\infty}
\frac{-ie^{i(\xi-1)t}}{\xi}\,\mathrm d\xi\\
&=e^{-it}\int_{-\infty}^{+\infty}
\frac{-ie^{i\xi t}}{\xi}\,\mathrm d\xi\\
&=e^{-it}
\begin{cases}
\pi,&t>0,\\
-\pi,&t<0.
\end{cases}
\end{aligned}
\]

Substitution gives
\[
\begin{aligned}
\mathcal F^{-1}[F(\omega)]
&=
\begin{cases}
\dfrac1{2\pi}(\pi e^{-it}+\pi e^{-it}),&t>0,\\[4pt]
\dfrac1{2\pi}(\pi e^{-it}-\pi e^{-it}),&t<0
\end{cases}\\
&=
\begin{cases}
e^{-it},&t>0,\\
0,&t<0
\end{cases}\\
&=u(t)e^{-it}.
\end{aligned}
\]

That is,
\[
u(t)e^{-it}
\ \xleftrightarrow{\mathcal F}\
\pi\delta(\omega+1)-\frac{i}{\omega+1}.
\]

Compare this with
\[
u(t)\ \xleftrightarrow{\mathcal F}\
\pi\delta(\omega)+\frac1{i\omega}
\]

The comparison shows a shift in the frequency domain.

\subsection{Shifting}
\[
f(t-t_0)
\ \xleftrightarrow{\mathcal F}\
e^{-it_0\omega}F(\omega).
\]

Alternatively,
\[
e^{i\omega_0t}f(t)
\ \xleftrightarrow{\mathcal F}\
F(\omega-\omega_0).
\]

\textbf{Example.} Starting from
\[
u(t)\ \xleftrightarrow{\mathcal F}\
\frac1{i\omega}+\pi\delta(\omega),
\]

Hence,
\[
u(t-c)\ \xleftrightarrow{\mathcal F}\
\left[\frac1{i\omega}+\pi\delta(\omega)\right]e^{-ic\omega}.
\]

This can be simplified to
\[
\frac{e^{-ic\omega}}{i\omega}+\pi\delta(\omega).
\]

Here we use the delta-function identity
\[
\delta(x-x_0)\varphi(x)
=\delta(x-x_0)\varphi(x_0).
\]

In addition,
\[
e^{-it}u(t)\ \xleftrightarrow{\mathcal F}\
\pi\delta(\omega+1)+\frac1{i(\omega+1)},
\]

and
\[
\mathcal F^{-1}\!\left[
\pi\delta(\omega-1)+\frac1{i(\omega-1)}
\right]
=u(t)e^{it}.
\]

\subsection{Scaling}

If \(f(t)\xleftrightarrow{\mathcal F}F(\omega)\), then
\[
f(at)\ \xleftrightarrow{\mathcal F}\
\frac1{|a|}F\!\left(\frac{\omega}{a}\right),
\qquad a\ne0.
\]

and
\[
\frac1{|a|}f\!\left(\frac{t}{a}\right)
\ \xleftrightarrow{\mathcal F}\
F(a\omega).
\]

\subsubsection*{Example: Find \(\mathcal F[u(5t-2)]\)}

\textbf{Solution (using the definition)}\par\nopagebreak[4]
\[
\mathcal F[u(5t-2)]
=\int_{-\infty}^{+\infty}
u(5t-2)e^{-i\omega t}\,\mathrm dt.
\]

Let \(5t-2=\tau\), so that \(t=(\tau+2)/5\). We obtain
\[
\begin{aligned}
\mathcal F[u(5t-2)]
&=\int_{-\infty}^{+\infty}
u(\tau)e^{-i\omega(\tau+2)/5}\frac15\,\mathrm d\tau\\
&=\frac15e^{-2i\omega/5}
\int_{-\infty}^{+\infty}
u(\tau)e^{-i(\omega/5)\tau}\,\mathrm d\tau\\
&=\frac15e^{-2i\omega/5}
\left[
\frac1{i(\omega/5)}+\pi\delta\!\left(\frac{\omega}{5}\right)
\right]\\
&=\frac{e^{-2i\omega/5}}{i\omega}
+\frac{\pi}{5}\delta\!\left(\frac{\omega}{5}\right).
\end{aligned}
\]

\subsection{Duality}

If
\[
f(t)\ \xleftrightarrow{\mathcal F}\ F(\omega),
\]

then
\[
F(t)\ \xleftrightarrow{\mathcal F}\ 2\pi f(-\omega).
\]

\textbf{Proof}\par\nopagebreak[4]
\[
F(\omega)
=\int_{-\infty}^{+\infty}f(t)e^{-i\omega t}\,\mathrm dt.
\]

A definite integral does not depend on the name of its dummy variable.
Interchanging \(\omega\) and \(t\) in the notation gives
\[
F(t)=\int_{-\infty}^{+\infty}f(\omega)e^{-it\omega}\,\mathrm d\omega.
\]

Let \(\omega=-\xi\).
\[
\begin{aligned}
F(t)
&=\int_{+\infty}^{-\infty}
f(-\xi)e^{it\xi}\,\mathrm d(-\xi)\\
&=\int_{-\infty}^{+\infty}
f(-\xi)e^{it\xi}\,\mathrm d\xi\\
&=\frac1{2\pi}\int_{-\infty}^{+\infty}
[2\pi f(-\omega)]e^{it\omega}\,\mathrm d\omega.
\end{aligned}
\]

Comparison with the inverse transform formula proves the result.

\textbf{Example.} Start from the transform pair for a two-sided
exponential pulse,
\[
e^{-\beta|t|}
\ \xleftrightarrow{\mathcal F}\
\frac{2\beta}{\omega^2+\beta^2},
\qquad \beta>0,
\]

we obtain
\[
\frac{2\beta}{t^2+\beta^2}
\ \xleftrightarrow{\mathcal F}\
2\pi e^{-\beta|-\omega|}
=2\pi e^{-\beta|\omega|}.
\]

Therefore,
\[
\frac1{t^2+\beta^2}
\ \xleftrightarrow{\mathcal F}\
\frac{\pi}{\beta}e^{-\beta|\omega|}.
\]

The two-sided exponential
pulse is symmetric about the vertical axis, reaches its maximum at the
origin, and decays in both directions.

\subsection{Differentiation}

\textbf{(1) Differentiation in the time domain}

If \(f(\pm\infty)=0\), then
\[
f'(t)\ \xleftrightarrow{\mathcal F}\
i\omega F(\omega).
\]

\textbf{Proof.} Integration by parts gives
\[
\begin{aligned}
\int_{-\infty}^{+\infty}f'(t)e^{-i\omega t}\,\mathrm dt
&=\left.f(t)e^{-i\omega t}\right|_{-\infty}^{+\infty}
-\int_{-\infty}^{+\infty}(-i\omega)f(t)e^{-i\omega t}\,\mathrm dt\\
&=0+i\omega\int_{-\infty}^{+\infty}
f(t)e^{-i\omega t}\,\mathrm dt\\
&=i\omega F(\omega).
\end{aligned}
\]

\textbf{(2) Differentiation in the frequency domain}
\[
-itf(t)\ \xleftrightarrow{\mathcal F}\ F'(\omega).
\]

\textbf{Proof}\par\nopagebreak[4]
\[
F(\omega)=\int_{-\infty}^{+\infty}
f(t)e^{-i\omega t}\,\mathrm dt,
\]
\[
\begin{aligned}
F'(\omega)
&=\int_{-\infty}^{+\infty}
\frac{\partial}{\partial\omega}
[f(t)e^{-i\omega t}]\,\mathrm dt\\
&=\int_{-\infty}^{+\infty}
f(t)(-it)e^{-i\omega t}\,\mathrm dt.
\end{aligned}
\]

\textbf{Generalization}\par\nopagebreak[4]
If
\[
f(\pm\infty)=f'(\pm\infty)=\cdots
=f^{(n-1)}(\pm\infty)=0,
\]

then
\[
f^{(n)}(t)\ \xleftrightarrow{\mathcal F}\
(i\omega)^nF(\omega).
\]

On the other hand,
\[
(-it)^nf(t)\ \xleftrightarrow{\mathcal F}\
F^{(n)}(\omega).
\]

From
\[
\mathcal F[(-it)^nf(t)]=F^{(n)}(\omega)
\]

we obtain
\[
\mathcal F[t^nf(t)]=i^nF^{(n)}(\omega).
\]

In particular,
\[
\begin{aligned}
\mathcal F[t^n\cdot1]
&=i^n[2\pi\delta(\omega)]^{(n)}\\
&=2\pi i^n\delta^{(n)}(\omega).
\end{aligned}
\]

That is,
\[
t^n\ \xleftrightarrow{\mathcal F}\
2\pi i^n\delta^{(n)}(\omega).
\]

\subsection{Examples of Generalized Fourier Transform Pairs}

All derivatives in the right-hand column are with respect to \(\omega\).

\begingroup
\small
\renewcommand{\arraystretch}{1.65}
\everymath{\displaystyle}
\begin{longtable}{@{}p{0.29\textwidth}p{0.65\textwidth}@{}}
\toprule
Function & Fourier transform \\
\midrule
\endhead
\bottomrule
\endlastfoot
\(u(t-c)\) &
\(\left[\frac1{i\omega}+\pi\delta(\omega)\right]e^{-ic\omega}\) \\
\(t\,u(t)\) &
\(i\frac{\mathrm d}{\mathrm d\omega}\left[\frac1{i\omega}+\pi\delta(\omega)\right]\) \\
\(t^nu(t)\) &
\(i^n\frac{\mathrm d^n}{\mathrm d\omega^n}\left[\frac1{i\omega}+\pi\delta(\omega)\right]\) \\
\(u(t)e^{iat}\) & \(\frac1{i(\omega-a)}+\pi\delta(\omega-a)\) \\
\(u(t)e^{iat}t^n\) &
\(i^n\frac{\mathrm d^n}{\mathrm d\omega^n}\left[\frac1{i(\omega-a)}+\pi\delta(\omega-a)\right]\) \\
\(\delta^{(n)}(t)\) & \((i\omega)^n\) \\
\(\delta^{(n)}(t-c)\) & \((i\omega)^ne^{-ic\omega}\) \\
\(t^ne^{iat}\) & \(2\pi i^n\delta^{(n)}(\omega-a)\) \\
\end{longtable}
\endgroup

The modulated transform pairs follow from the frequency shift
\(\omega\mapsto\omega-a\).

\subsection{Integration}

\textbf{(1) Integration in the time domain}

If
\[
\int_{-\infty}^{+\infty}f(t)\,\mathrm dt=0,
\]

then
\[
\int_{-\infty}^{t}f(\tau)\,\mathrm d\tau
\ \xleftrightarrow{\mathcal F}\
\frac{F(\omega)}{i\omega}.
\]

\textbf{(2) Integration in the frequency domain}

If \(F\) is sufficiently regular and decays at infinity, with
\[
\int_{-\infty}^{+\infty}F(\xi)\,\mathrm d\xi=0,
\]
then, under the corresponding transform-existence conditions,
\[
\frac{f(t)}{-it}
\ \xleftrightarrow{\mathcal F}\
\int_{-\infty}^{\omega}F(\xi)\,\mathrm d\xi.
\]

\textbf{Proof of (2)}\par\nopagebreak[4]
Let
\[
G(\omega)=\int_{-\infty}^{\omega}F(\xi)\,\mathrm d\xi,
\]

Differentiation gives \(F(\omega)=G'(\omega)\). By the differentiation
property,
\[
\begin{aligned}
\mathcal F^{-1}[F(\omega)]
&=\mathcal F^{-1}[G'(\omega)]\\
&=-it\,\mathcal F^{-1}[G(\omega)].
\end{aligned}
\]

Hence,
\[
\mathcal F^{-1}[G(\omega)]
=\frac{\mathcal F^{-1}[F(\omega)]}{-it}
=\frac{f(t)}{-it}.
\]

That is,
\[
\mathcal F^{-1}\!\left[
\int_{-\infty}^{\omega}F(\xi)\,\mathrm d\xi
\right]
=\frac{\mathcal F^{-1}[F]}{-it}.
\]

\textbf{Proof of (1)}\par\nopagebreak[4]
Let
\[
g(t)=\int_{-\infty}^{t}f(\tau)\,\mathrm d\tau,
\]

Then \(f(t)=g'(t)\). If \(g(\pm\infty)=0\), the differentiation property
gives
\[
\mathcal F[f(t)]
=\mathcal F[g'(t)]
=i\omega\mathcal F[g(t)].
\]

It follows that
\[
\mathcal F[g(t)]=\frac{F(\omega)}{i\omega}.
\]

If
\[
\int_{-\infty}^{+\infty}f(t)\,\mathrm dt\ne0,
\]

then
\[
\mathcal F\!\left[
\int_{-\infty}^{t}f(\tau)\,\mathrm d\tau
\right]
=\frac{F(\omega)}{i\omega}
+\pi\delta(\omega)F(0).
\]
\end{document}
